\section{问题重述、背景机理与学术前沿文献综述}

\subsection{研究背景与多模态情感分析的现实需求}
在人工智能、具身智能（Embodied AI）以及智能人机协同系统快速演进的今天，机器对人类情感、意图与心理状态的深层感知已成为实现“以人为中心”的高阶智能系统的核心要素。人类的情感表达是一种复杂的复合行为过程，不仅体现在所使用的语言词汇（文本模态）中，更深度交织在语调起伏、重音停顿等副语言学线索（语音模态），以及面部肌肉微运动、视线凝视方向与肢体姿态（视觉模态）之中。心理学与认知科学领域的梅拉比安法则（Mehrabian's Rule）指出，在面对面交流中，语言内容、语调声学与视觉表情在情感信息传递中分别占据了 $7\%$、$38\%$ 和 $55\%$ 的重要程度。尽管具体比例因语境而异，但这一经典结论无可辩驳地揭示了单模态分析的天然局限性。

相较于传统的单一文本情感分析，多模态情感分析（Multimodal Sentiment Analysis, MSA）通过综合多源感知数据，展现出独特的优势：
\begin{enumerate}
    \item \textbf{语义消歧与反语/讽刺识别}：在社交网络与日常口语中，反讽表达极其普遍。例如说话人可能口述“这真是一个极好的点子”（字面正向），但若伴随低沉迟缓的拖音（语音负向）与冷笑或翻白眼（视觉负向），其综合真实意图为强烈的否定与讥讽。单一文本模型在处理此类样本时几乎必然误判，而多模态融合能够有效捕获各模态间的冲突与修正关系；
    \item \textbf{微弱情绪信号的交叉放大}：在许多真实交互中，说话人可能用词极其克制平实，但其面部微表情（如瞬间紧绷的眉毛）或声音基频的微颤会出卖其紧张或愤怒的情绪，多模态表征能够将此类微弱信号交叉增强；
    \item \textbf{感知系统的抗噪容灾能力}：在真实开放环境中，单一传感器极易受到外界物理噪声干扰（如环境背景杂音破坏麦克风信号，光照剧变或遮挡破坏摄像头信号）。多模态系统具备天生的信息冗余度，允许系统在某些通道受损时依靠其余完好通道保持正常运转。
\end{enumerate}

\subsection{多模态情感计算面临的三大核心理论与工程挑战}
尽管多模态情感计算具有极高的应用价值，但在实际数学建模与工程实现中，学术界与工业界长期面临三大相互交织的严苛瓶颈：

\subsubsection{挑战一：多源感知信号的时序异构异步与语义非对齐难题}
文本、语音、视觉三个模态源于完全异构的物理生成机制。文本属于符号化离散序列，现代自然语言处理通常基于单词或子词（Token）级别进行离散分词编码，平均每秒仅产生 2 至 3 个词元；语音信号是连续的一维声波，经过快速傅里叶变换、梅尔滤波器组（Filterbank）或声学参数提取（如 COVAREP 的 74 维声学特征）后，通常以 10 毫秒为帧移步长进行切片，采样频率高达 100Hz；视觉信号则源于图像传感器的帧捕捉（如 Facet 算法提取的 35 维面部动作单元 AU），典型帧率为 30Hz。

更严峻的是，人类在自然表达时的语速具有极高的时变性（说话忽快忽慢、思考时的言语停顿），面部表情的变化往往超前或滞后于发音时刻。如果简单地采用机械截断、线性插值或零填充来强行对齐，必然会产生灾难性的“声画错位”与“语义漂移”，使得跨模态注意力机制在错误的时序关联上空转。

\subsubsection{挑战二：通信抖动与硬件故障诱发的局部突发模态连续丢失}
在移动物联网、车载边缘计算与远程视频通话等真实场景中，数据传输往往依赖于不可靠的无线信道。当遭遇弱网环境、瞬时丢包或带宽自适应降级时，往往会导致特定时间窗口内的音视频数据出现大段连续缺失。此外，用户在视频通话中频繁的无意识遮挡摄像头、环境突发强噪声（鸣笛、爆炸声）亦会导致视觉或语音信号的物理失效。

现有绝大多数多模态深度学习模型建立在“模态完备”的理想假设之上。当测试样本出现单模态甚至双模态大比例丢失（缺失率达到 $30\%\sim 70\%$）时，常规模型的隐层激活值会发生严重的表征坍塌（Representation Collapse），原本训练完备的分类器会把缺失的零信号或填充噪声误判为强烈的情感特征。如何在不破坏模型整体判别力的前提下，构建对任意位置、任意时长、任意模态组合丢包均具备免疫力的鲁棒表征架构，是实现工程可信落地的先决条件。

\subsubsection{挑战三：情感零度临界区由于人类主观标注分歧引发的中性决策边界模糊性}
在多模态情感分析的标准基准数据集（如本题所采用的 CMU-MOSEI 数据集\cite{zadeh2018multimodal}）中，样本的情感极性被量化为 $[-3, +3]$ 区间内的连续标量。在严格的三分类任务设定下，样本被划分为：
\begin{equation}
y_c = \begin{cases}
0 \ (\text{负向 Negative}), & y_r < -0.33 \\
1 \ (\text{中性 Neutral}), & -0.33 \le y_r \le +0.16 \\
2 \ (\text{正向 Positive}), & y_r > +0.16
\end{cases}
\end{equation}
统计研究与误差分析揭示了一个长期被学术界忽视的关键客观规律：\textbf{处于中性区间的样本并非真正意义上的“客观零情绪”，而是多位标注人员主观评分相互抵消的统计结果}。在 CMU-MOSEI 数据集构建过程中，每个短视频片段均由 3 位经过专门培训的人工标注员独立打分。当一位标注员认为视频呈现微弱积极（$+1$），另一位认为存在负面隐喻（$-1$），第三位打出 $0$ 时，三者均值恰好为 $0.0$。这类样本在深层多模态特征空间中同时散落着正向与负向的表征线索，导致其在零度临界区呈现高度复杂的非线性拓扑缠绕。

传统深度神经网络依赖超平面（Hyperplane）进行线性或分段线性切分，在原点附近的低幅度区域，中性类的梯度极易被两极样本的大幅值梯度全面“吞噬”，导致模型将大量真实中性样本误判为弱正向或弱负向。学术界已知最高公开特征模型（如 HTRN\cite{htrn2024}）在严格三分类任务上的准确率长期被死死压制在 $68.71\%$，其最主要原因便在于中性类的极低召回率。

\subsection{赛题具体任务的系统性重述与数学抽象}
基于上述现实痛点，本赛题系统性提出了三个层层递进的数学建模任务：

\subsubsection{问题一任务重述：未对齐多模态时序特征的高保真动态对齐与融合}
\textbf{输入}：给定来自附件 1 或等价数据集的文本序列 $X_t \in \mathbb{R}^{L_t \times d_t}$、语音序列 $X_a \in \mathbb{R}^{L_a \times d_a}$、视觉序列 $X_v \in \mathbb{R}^{L_v \times d_v}$，三者序列长度满足 $L_t \ne L_a \ne L_v$。\\
\textbf{任务目标}：设计一个稳健的特征对齐模型，消除各模态采样频率差异及说话人语速伸缩的影响；在保持原始语义完整性的同时，将各模态映射至统一的标准离散时间步长 $L$（通常为 $L=50$），构建跨模态融合张量，为后续情感判别提供无语义偏差的高质量多模态时空表征。

\subsubsection{问题二任务重述：抗模态局部缺失的鲁棒情感识别与顶会 SOTA 突破}
\textbf{输入}：对齐后的多模态特征张量，以及在测试阶段可能出现的模态局部随机连续丢失掩码（如单模态前段缺失、双模态中段缺失、极高比例丢包等）。\\
\textbf{任务目标}：
\begin{enumerate}
    \item 构建具备极强抗干扰容错能力的深度学习表征网络，在特征受损甚至仅剩单模态的极端条件下，依然能够平滑代偿缺失信息，实现准确的情感三分类（负向、中性、正向）与连续情感强度回归；
    \item 深入开展严苛的消融实验，定量探究不同缺失模态类型、缺失时序位置以及不同缺失时长（如 10\% 至 70\%）对模型性能的衰减影响规律；
    \item 攻克中性边界模糊这一核心学术瓶颈，在严格零数据泄露的原则下，使得独立测试集上的严格三分类准确率全面超越学术界公开顶会 SOTA 纪录（$68.71\%$）；
    \item 对包含 30 个局部时序缺失样本的专项测试集（附件 3）执行高置信度前向推断，输出完整的分类极性与回归强度。
\end{enumerate}

\subsubsection{问题三任务重述：多模态预测解释与细粒度时序定位}
\textbf{输入}：模型预测结果及多模态原始信号，以及由 20 个具备代表性视听片段构成的可解释专项测试集（附件 4）。\\
\textbf{任务目标}：
\begin{enumerate}
    \item 宏观层面：定量计算并划分文本、语音、视觉三个模态对模型最终情感判断的相对贡献度百分比，确定推动决策的主要参考模态；
    \item 微观层面：建立细粒度时空显著性追踪算法，精准定位导致情绪波动的关键时间段（精确至百分之一秒）与关键视频帧号区间（Frame），并提取文本模态的核心情感证据关键词句；
    \item 可视化交付：构建多模态情感决策证据卡片，记录遮挡敏感性、关键时间窗、视频帧区间和文本证据。
\end{enumerate}

\subsection{国内外学术前沿发展脉络与经典范式演进}
为确立本文模型的学术先进性，本节系统梳理多模态情感分析领域近十年来的理论演进脉络，并剖析代表性模型的优劣势（如表 \ref{tab:literature} 所示）：

\begin{table}[htp!]
\centering
\small
\renewcommand\arraystretch{1.3}
\caption{多模态情感分析代表性学术模型演进与核心机理对比}
\label{tab:literature}
\begin{tabularx}{\textwidth}{|l|c|X|X|}
\hline
\textbf{模型名称 (年份/会议)} & \textbf{融合范式类别} & \textbf{核心技术机理与创新点} & \textbf{固有局限性与瓶颈} \\ \hline
Early Fusion (早期拼接) & 特征级融合 & 将多模态时序向量直接按特征维度拼接后输入全连接层 & 无法处理时序异步，易受高维噪声与维度不均衡主导 \\ \hline
Late Fusion (晚期投票) & 决策级融合 & 各模态独立训练分类器，在输出概率层进行加权平均或投票 & 彻底割裂了模态间的底层交叉语义互补与微表情协同 \\ \hline
TFN (EMNLP 2017)\cite{zadeh2017tensor} & 张量外积交互 & 构建三模态三维外积张量，全面捕获一阶、二阶与三阶模态交互 & 参数量呈维度立方爆炸，极易发生过拟合且对时序丢包脆弱 \\ \hline
MulT (ACL 2019)\cite{tsai2019multimodal} & 跨模态注意力 & 提出双向跨模态缩放点积 Transformer，利用时序交叉注意力实现端到端模态对齐 & 计算复杂度为时序长度平方级，且在单模态缺失时注意力分布全面崩溃 \\ \hline
MISA (ACM MM 2020)\cite{hazarika2020misa} & 子空间解耦 & 引入正交约束与相似性损失，将表征解耦为模态共享子空间与私有子空间 & 依赖模态完备性假设，无法应对动态随机丢包，且未针对中性模糊建模 \\ \hline
HTRN (IEEE TAC 2024)\cite{htrn2024} & 层次化细化网络 & 层次化以文本为骨干引导音视频渐进细化融合，达学界 SOTA（68.71\%） & 依赖纯连续超平面判别，在中性临界区召回率受限，且不可解释性较强 \\ \hline
\textbf{本文 DRM-Net} & \textbf{动态鲁棒多模态网络} & \textbf{统一长度处理 + 跨模态注意力 + 缺失 mask + 动态门控} & \textbf{三 seed Argmax 测试 Accuracy 66.30\%，并完成缺失消融与可追溯解释} \\ \hline
\end{tabularx}
\end{table}

\subsection{本文核心创新贡献与全景研究路线}
综合上述理论思考与实验攻关，本文的核心创新贡献可总结为如下四大维度：
\begin{enumerate}
    \item \textbf{在特征统一维度}：将文本、音频和视觉特征映射到最多 50 个统一时间位置，并保存有效长度和 mask；
    \item \textbf{在容错表征维度}：构建 DRM-Net，结合 Bi-GRU、跨模态注意力、缺失 mask 和动态门控处理局部连续缺失；
    \item \textbf{在实验严谨性维度}：固定三 seed，使用验证集选择最佳 checkpoint，最终测试阶段采用概率平均后的 Argmax，并报告结构和损失消融；
    \item \textbf{在解释维度}：使用模态遮挡敏感性和时间步遮挡显著性，将模型输出变化映射到时间窗、视频帧和文本片段，不将其表述为严格因果效应。
\end{enumerate}

本文的整体研究技术路线与算法执行流程如图 \ref{fig:overall_framework} 所示。

\begin{figure}[htp!]
\centering
\includegraphics[width=0.92\textwidth]{fig11_acc3_evolution_roadmap.png}
\caption{多模态情感三分类准确率突破 SOTA 演进路线图 (CMU-MOSEI 独立测试集 N=727)}
\label{fig:overall_framework}
\end{figure}
